\subsection{AI Systems that Adapt to Novel Situations}
\label{sec:6.2.1}
The social and political situations this book cares about are open-ended, so a system that holds its commitments only inside its training envelope will not hold them for long. The technical machinery for adapting is covered where it belongs: meta-learning and transfer learning at sections 4.1.3 and 4.2.4, lifelong learning at section 4.1.2, and active learning — a system routing its least confident cases to a person instead of resolving them alone — is the one addition worth naming here, because it is the mechanism by which a system can be uncertain in public.

What matters at this point in the argument is a distinction the word adaptability hides. A system has to adapt to a changing environment without its commitments adapting along with it, and these are not two settings of one dial. Erosion by adaptation is how an ethical constraint disappears without anyone deciding to remove it: each individual update is defensible against the case in front of it, and the sequence goes somewhere nobody chose. Chapter 3's floor is the structural answer, and section 3.3 is where the cost of holding one against the party doing the updating gets counted.

Two consequences follow that the rest of the book relies on. A system that keeps changing is harder to audit than a static one, so transparency requirements have to attach to the update process and not only to a released version. And high adaptability is an attack surface: a system built to change in response to its environment can be moved by an adversary who arranges the environment, which is section 10.3.3's problem in its ordinary, non-state form.
